\chapter[Memory across several reporting periods]{Memory across reporting periods: taxes, vintages, arrears, and future rights}
\label{sim:memory-case}

Two companies with the same current statements can differ because their past creates different future rights and obligations. A simulation that keeps only current totals forgets this memory. This chapter follows several small examples over more than one period. The examples use stipulated accounting and tax rules solely to expose state requirements. They are not claims about a particular jurisdiction's current law.

\section{A loss carryforward}
Assume a teaching tax system with rate 25 per cent, no refund for current losses, unlimited use of carried losses against future positive taxable income and no expiry. Let $L_t$ be unused losses entering period $t$ and $T_t$ taxable income before using them. For positive $T_t$, used losses are $U_t=\min(L_t,T_t)$. Cash tax and next loss balance are
\begin{align}
 \operatorname{Tax}_t&=0.25\max(T_t-L_t,0),\\
 L_{t+1}&=\max(L_t-\max(T_t,0),0)+\max(-T_t,0).
\end{align}
The second equation covers both profit and loss periods under these simplified rules. A positive income reduces available losses; a negative income adds its magnitude.

Company A enters with losses 40 and company B with none. Both have the same current cash, assets and book equity because the example assumes no recognized deferred tax asset. Next period taxable income is 30. A pays no cash tax and retains losses 10. B pays 7.5. Current statements alone do not identify that difference.

In the following period taxable income is again 30. A uses its remaining 10 and pays 5; B pays 7.5. Over two periods, tax is 5 versus 15. Omitting loss memory overstates A's tax. Treating the original loss as an immediate refund would instead overstate its initial cash. Both errors can occur in simple effective-tax-rate forecasts.

\section{Limits and expiry change the recurrence}
Now impose a rule that losses can offset at most 50 per cent of positive taxable income in a period. With $L_0=40$ and $T_0=30$, usable losses are 15, tax is 3.75 and remaining losses are 25. The earlier recurrence is no longer correct.

An expiry date requires vintage-specific loss balances. A single total 40 cannot determine which amount expires next year. Two companies with the same total can have different future tax capacity. The state may need a vector by expiry cohort, jurisdiction and restriction.

A change of ownership can affect utilization under some actual systems, but no such rule is assumed here without authority. The point is structural: tax rights depend on facts beyond a total balance. The model should admit the applicable rule and preserve the state it requires.

\section{Deferred tax is a measurement, not the loss pool}
A deferred tax asset may be recognized only under specified accounting criteria. Its carrying amount can differ from the mechanical tax value of losses. A valuation allowance or recognition restriction can change book equity without changing the legal loss pool.

Suppose the legal losses remain 40 while a recognized deferred tax asset is reduced from 10 to 4 under the chosen assessment. The reduction is not a cash tax payment. The future tax computation still needs the legal loss balance and applicable utilization rules.

Conversely, recognizing an asset of 10 does not create 10 of spendable cash. A liquidity scenario should use actual tax payment or refund dates. The accounting measurement and economic tax rights are linked but distinct states.

\section{Asset vintages with the same carrying amount}
Consider two machines under a simple straight-line model with zero residual value. Machine A originally cost 100 and has two of ten years remaining, so its carrying amount is 20 and annual depreciation 10. Machine B originally cost 40 and has one of two years remaining, also carrying 20 but annual depreciation 20. Equal current carrying amount does not imply equal next-year expense.

If the reported current depreciation is also constrained to match, combinations of several machines can still create ambiguity. For example, a portfolio can mix assets with different remaining lives whose total charge is the same. The next-period charge changes when a cohort reaches the end of life.

The state should retain acquisition or remaining-life cohorts where depreciation, maintenance or replacement depends on them. A single average life may be adequate for a smooth aggregate forecast but can miss a replacement wall. The approximation should be checked against the horizon and decision.

\section{Replacement and operating capacity}
An asset can be fully depreciated while still productive, or retain book value while becoming unusable. Depreciation is not a direct physical failure process. Replacing capital expenditure with depreciation assumes a relationship that must be justified separately.

Suppose A's machine can operate for two more years without replacement and B's fails after one year unless 30 is spent. Both enter with carrying amount 20. A financing shock in year two affects B's production but not A's in the same way. A model using only carrying amount cannot distinguish them.

A richer state includes capacity, maintenance, failure or retirement rule and replacement lead time. These can be uncertain. The history family may retain alternative vintage and condition profiles compatible with the same accounting data. Their scenario outcomes remain conditional rather than being collapsed into one average machine.

\section{Arrears are not current-period expense again}
Suppose a loan coupon of 5 accrues in period one but is not paid. Under the illustrative accrual treatment, interest expense and a payable increase by 5. In period two, paying the arrears reduces cash and the payable. It does not recognize the same interest expense again.

If penalty interest accrues on arrears, that is an additional contractual event with its own base, rate and dates. The state must retain unpaid amounts and payment allocation rules. A payment may be applied first to fees, interest or principal according to the admitted contract.

A model that stores only total debt can lose this allocation. Paying 5 might reduce principal when it should settle interest, altering future coupon calculations. The resulting statements may still balance, but the contract state becomes wrong. Multi-period replay exposes the error.

\section{A payment allocation example}
At a date, principal is 100, accrued interest 5 and fee arrears 2. A payment of 6 is made. Under the stipulated order fees first, then interest, then principal, fees become zero, interest becomes 1 and principal remains 100.

Under principal-first allocation, principal becomes 94 while interest and fees remain 7. Total obligation falls by 6 in both cases, but future interest and default status can differ. A total-debt bridge cannot distinguish the two. The allocation rule is part of the contract.

If the rule is unknown, both allocation histories may remain possible. A scenario should not choose the order that minimizes future cost unless management actually has that right. Creditor-controlled allocation cannot be optimized by the borrower without permission in the model.

\section{Cure periods and clocks}
A breach can start a cure period measured in business or calendar days. The state needs the trigger date, deadline, required remedy and whether cure has occurred. Merely retaining a Boolean breach flag loses the deadline.

Suppose a payment is missed on day 1 and can be cured within three calendar days, with the deadline defined at the end of day 4 under the stipulated convention. A receipt on day 4 can cure if settlement precedes the cutoff; one on day 5 cannot. An annual or monthly timestamp cannot resolve this distinction.

The actual contract may count days differently, include notices or exclude holidays. The simulator should not infer a convention from a generic date library. Exact boundary fixtures should accompany the admitted timing rule.

\section{Options survive quiet periods}
An option can remain out of the money for many reporting periods and then become decisive under a shock. A call right, put right, conversion feature or renewal option should not disappear because it had zero exercise activity historically.

Two loans can have identical coupon and principal histories while one permits early redemption at par and the other imposes a premium. A rate decline creates different refinancing choices. Current book debt and past interest do not identify the option.

A state summary can omit the option only for a query where its exercise cannot affect outputs over the horizon. That claim may rely on bounds showing the option never becomes relevant. Once the scenario or horizon changes, the summary must be reconsidered.

\section{A refinancing option value example}
Suppose a company owes a fixed payment of 110 next year. It can refinance today at a cost of 102 if it has a call right, and current cash is sufficient. Ignoring discounting for this one-period comparison, the right can save 8. Without the right, the old obligation remains.

If exercising requires a fee of 9, the saving disappears and the policy should not exercise under a cost-minimizing objective. If the fee is uncertain between 0 and 12, the exercise decision is conditional. A model that always redeems at par gives every borrower an option it may not possess.

The example does not price the option probabilistically. It shows a state-contingent decision consequence. Pricing would require a distribution or valuation operator over future rates and policies, along with the contractual exercise schedule.

\section{Estimates carry information history}
An impairment estimate can depend on forecasts and discount rates available at the reporting date. A later recovery in value does not mean the earlier estimate was necessarily an error. A restatement that corrects a mistaken fact is different from an estimate revised because new information arrived.

A historical simulation should preserve the information used for each estimate. If it recomputes every past estimate using today's knowledge, it changes the as-reported history and can introduce look-ahead into validation. A current restated view can be useful, but should be named separately.

The economic asset state and accounting estimate can therefore have different version histories. A new estimate changes measurement without necessarily changing physical assets or cash. A model should track the relation rather than force every valuation change into an economic transaction.

\section{Inventory age and obsolescence}
Inventory with the same quantity and carrying value can have different remaining shelf life. A demand slowdown can make one cohort obsolete while another remains saleable. A pooled average age can conceal the threshold.

Suppose ten units expire next month and ten otherwise identical units expire next year. Both pools are carried at 50. If no sales occur for two months, the first may lose all saleable value under the stipulated physical rule while the second remains available. Historical cost equality does not preserve the response.

A cohort model can retain quantities by expiry and location without individual serial numbers when units within a cohort are interchangeable for the query. This is a justified aggregation opportunity. If a recall affects a particular batch, the cohort key must include batch identity or be refined.

\section{Customer behavior has memory}
A collection process can depend on invoice age, previous delinquency, dispute status and payment promises. A single receivable total or average collection ratio may omit these predictors. Two companies with equal receivables can have different next-month receipts.

A Markov model can be appropriate if the chosen state captures the relevant dependence under its assumptions. Calling the process Markov does not make a deficient state sufficient. If repayment depends on the duration of arrears, include age or a suitable duration state.

An empirical model should test whether the proposed summary predicts the relevant transitions and whether stress changes that relation. A purely contractual due-date model answers a different question from a behavioral collection model. Both can be retained with separate roles.

\section{Horizon truncation}
A policy can look safe over one year by issuing debt due just after the horizon. The closing cash is positive, but the future obligation remains. A terminal report should show that maturity and avoid treating its proceeds as permanent wealth.

Suppose a company borrows 20 at year-end and repays 22 one month later. A one-year cash objective rewards the 20 inflow unless terminal debt is valued or constrained. A policy optimizer can exploit this omission even if the event mechanics are correct.

Possible remedies include extending the horizon, adding a justified continuation value, requiring terminal funding coverage or reporting residual obligations separately. Each changes the decision problem in a stated way. No terminal rule is universally correct for all company questions.

\section{A sufficient state is query-relative}
For a one-day payment query, asset age may be irrelevant while settlement dates are essential. For a five-year capital plan, vintages and replacement needs matter. For tax cash, legal loss pools matter; for book equity, deferred tax measurement also matters.

The underlying representation should preserve these quantities where supported, while the active execution state can be reduced for the query. A sufficiency argument names the allowed scenarios, policies, outputs and horizon. It is not a permanent declaration that discarded details never matter.

This is how compact simulation can coexist with rich historical evidence. The system need not replay every old event at every step, but it must retain the memory that determines future events and the provenance needed to reconstruct it after revisions.

\section{Implementation tests across periods}
A loss-pool test should distinguish cash tax from deferred tax measurement and preserve expiry cohorts when required. A vintage test should produce different replacement paths from equal carrying amounts. An arrears test should prevent double expense and apply payments in the stipulated order.

A cure-clock test should distinguish receipts before and after the deadline. An option test should preserve unexercised rights through quiet periods. A horizon test should retain obligations beyond the final reporting date. Each case targets a specific missing state variable.

The tests should also verify legitimate aggregation. Identical loans with the same terms and dates can sometimes be combined additively. Identical inventory units can be pooled by cohort. The goal is sufficient memory, not maximal detail for its own sake.
